\documentclass[11pt]{article}

\usepackage{extras} % Se extras.sty

\begin{document}
\begin{titlepage}
\begin{center}

{\Large\bfseries TSEA56 - Kandidatprojekt i elektronik \\ LIPS Teknisk Dokumentation}

\vspace{5em}

Version 1.0

\vspace{5em}
Grupp 4 \\
\begin{tabular}{rl}
Hynén Ulfsjöö, Olle&\verb+ollul666+
\\
Wasteson, Emil&\verb+emiwa068+
\\
Tronje, Elena&\verb+eletr654+
\\
Gustafsson, Lovisa&\verb+lovgu777+
\\
Inge, Zimon&\verb+zimin415+
\\
Strömberg, Isak&\verb+isast763+
\\
\end{tabular}

\vspace{5em}
\today

\vspace{16em}
Status
\begin{longtable}{|l|l|l|} \hline

Granskad & - & - \\ \hline
Godkänd & - & - \\ \hline
 
\end{longtable}

\end{center}
\end{titlepage}

\pagebreak
\begin{center}

\section*{PROJEKTIDENTITET}
2016/VT, Undsättningsrobot Gr. 4

Linköpings tekniska högskola, ISY
\vspace{5em}
%\begin{center}

\begin{tabular}{|l|l|l|l|} \hline
\textbf{Namn} & \textbf{Ansvar} & \textbf{Telefon} & \textbf{E-post}  \\ \hline 
Isak Strömberg (IS) & Projektledare & 073-980 38 50 & isast763@student.liu.se \\ \hline
Olle Hynén Ulfsjöö (OHU)& Dokumentansvarig & 070-072 91 84 & ollul666@student.liu.se \\ \hline
Emil Wasteson (EW) & Hårdvaruansvarig & 076-836 61 66 & emiwa068@student.liu.se \\ \hline
Elena Tronje (ET) & Mjukvaruansvarig & 072-276 92 93 & eletr654@student.liu.se \\ \hline
Zimon Inge (ZI)& Testansvarig & 070-171 35 18 & zimin415@student.liu.se \\ \hline
Lovisa Gustafsson (LG) & Leveransansvarig & 070-210 32 53 & lovgu777@student.liu.se \\ \hline
\end{tabular}

%\end{center}

E-postlista för hela gruppen: isast763@student.liu.se

\vspace{5em}
Kund: ISY, Linköpings universitet 
tel: 013-28 10 00, fax: 013-13 92 82 \\
Kontaktperson hos kund: Mattias Krysander \\
tel: 013-28 21 98, e-post: matkr@isy.liu.se \\

\vspace{5em}
Kursansvarig:  Tomas Svensson\\
tel: 013-28 13 68, e-post: tomass@isy.liu.se \\
Handledare: Peter Johansson \\
tel: 013-28 13 45, e-post: peter.a.johansson@liu.se
\end{center}
\pagebreak

\tableofcontents

\pagebreak

\section*{Dokumenthistorik}
\begin{table}[h]
\begin{tabular}{|l|l|l|l|l|} \hline

\textbf{Version} & \textbf{Datum} & \textbf{Utförda förändringar} & \textbf{Utförda av} & \textbf{Granskad} \\ \hline
0.1 & - &  Första utkastet & Grupp 4 & - \\ \hline
\end{tabular}
\end{table}

\pagebreak
\pagenumbering{arabic}

\begin{flushleft}
\section{Inledning}
\textit{Bakgrund och syfte}
Det genomförda projektet har gått ut på att konstruera en undsättningsrobot. Detta dokument syftar till att användas som underlag till personer som vill bygga en liknande robot. Innehållet består av detaljerade beskrivningar av systemets olika delar och dess mjuk- och hårdvarukomponenter samt implementation.

\section{Produkten}
\textit{En bild på produkten och en beskrivning av hur den fungerar. Beskriv vad den används till.}
Produkten som har konstruerats är en undsättningsrobot bestående av fyra moduler varav tre är, tillsammans med sensorer och en gripklo, integrerade på ett gemensamt chassi enligt figur \ref{overview}. Den fjärde modulen är en dator som visar lämplig information. ÄNDRA TILL 1284P OCH EV GP2D120

\begin{figure}[!htbp]
\centering
\noindent\resizebox{\linewidth}{!}{
	
\documentclass[border=10px]{standalone}
\usepackage{tikz}
\usetikzlibrary{patterns}
\usetikzlibrary{shapes.arrows}
\usepackage{amssymb}
\usetikzlibrary{calc}
\usepackage{verbatim}
\usetikzlibrary{decorations.pathmorphing}

\begin{document}
	
\begin{tikzpicture}[scale=1,rotate=0]
		
	%Base
	\draw[thick, draw=black, fill=gray!10] (0,0) rectangle (6,10);

	%Wheels
	\draw[thick, pattern=north west lines, pattern color=black] (-.5,1) 		rectangle (0,2.5);
	\draw[thick, pattern=north west lines, pattern color=black] (-.5,7.5) 	rectangle (0,9);
	\draw[thick, pattern=north west lines, pattern color=black] (6,1) 		rectangle (6.5,2.5);
	\draw[thick, pattern=north west lines, pattern color=black] (6,7.5) 		rectangle (6.5,9);
	
	%Sensors
	\draw[thick, draw=black, fill=white]		(-.25,.25) 		rectangle 	(.5,.75);
	\draw[thick, draw=black, fill=white] 	(-.25,9.25) 		rectangle 	(.5,9.75);
	\draw[thick, draw=black, fill=white] 	(5.5,.25) 		rectangle 	(6.25,.75);
	\draw[thick, draw=black, fill=white] 	(5.5,9.25) 		rectangle 	(6.25,9.75);
	\draw[thick, draw=black, fill=white] 	(1.5,10.25) 		rectangle 	(4.5,9.5);
	\draw[thick, draw=black]				 	(2.5,9.5)			--		  	(2.5,10.25);
	\draw[thick, draw=black]					(3.5,9.5) --++ (0,.75);
	
	%Gripklo
	\draw[thick, draw=black, <-] (3,10.25) --++ (0,.72) node[above] {Gripklo};
	
	%Lidar lite
	\draw[thick, draw=black, <-] 			(1.5,10.25)  		--+ 		  	(135:1) node[above, align=center] {\verb+LIDAR-Lite v2+ \\ med servo};
	%Identifierare av nödställd
	\draw[thick, draw=black, <-]				(4.5,10.25)		--+			(45:1)  node[above] {\verb+IRM-8601-S+};
	
	%IR sensorer
	\draw[thick, draw=black, <-]				(-.25, 9.75) 	--+			(135:1) node[above] {\verb+GP2D120+};
	\draw[thick, draw=black, <-]				(-.25, .25) 		--+			(225:1) node[below] {\verb+GP2D120+};
	\draw[thick, draw=black, <-]				(6.25, 9.75) 	--+			(45:1) node[above] {\verb+GP2D120+};
	\draw[thick, draw=black, <-]				(6.25, .25) 		--+			(315:1) node[below] {\verb+GP2D120+};
	
	\draw[thick, draw=black, fill=white]					(2.5,4.5)		rectangle	(3.5,5.5);
	
	%Gyro
	\draw[thick, draw=black, <-] (3.5,5) --++ (.25,0) node[rotate=270, above] {\verb+MLX90609+};
	
	%Arrows and text
	%\draw[thick, ->]  (3,11) node[left, align=center] {\verb+LIDAR-Lite v2+ \\ + detektor av nödställd} -- (3,10.25);
	%\draw[thick, <->] (0.5,0.5)  --  (5.5,0.5) node[left=-14pt,midway, fill=gray!10] {\verb+GP2D120+};
	%\draw[thick, <->] (0.5,9.5) -- (3,9) node[right=-14pt,fill=gray!10] {\verb+GP2D120+} -- (5.5,9.5);
	%\draw[thick, ->] (4.5,5) node[above] {\verb+MPU-6500+} -- (3.5,5);
	
	%Sensormodul
	\node (sensormodul) at (3,6.75) [thick, draw=black, minimum width=3cm, minimum height=1.5cm, align=center, fill=white] {Sensormodul \\ \verb+ATmega1284p+};
	
	%Huvudmodul
	\node (huvudmodul) at (3,3.25) [thick, draw=black, minimum width=3cm, minimum height=1.5cm, align=center, fill=white] {Huvudmodul \\ \verb+ATmega1284p+};
	
	%Styrmodul
	\node (styrmodul) at (3,1.25) [thick, draw=black, minimum width=3cm, minimum height=1.5cm, align=center, fill=white] {Styrmodul \\ \verb+ATmega1284p+};
	
	%I2C-buss
	\draw[thick, draw=teal] ($(sensormodul.north east) + (.25,0)$) -- ($(styrmodul.south east) + (.25,0)$);
	\draw[thick, draw=teal] ($(sensormodul.north east) + (.5,0)$) -- ($(styrmodul.south east) + (.5,0)$);
	\node (i2c) at ($(sensormodul.north east) + (0.375,.25)$) [text=teal] {\small I\textsuperscript{2}C};
	
	\draw[thick, draw=teal, fill=teal] ($(sensormodul.south east) + (0,.25)$) --+ (.25,0) circle [radius=.05cm];
	\draw[thick, draw=teal, fill=teal] ($(sensormodul.south east) + (0,.5)$) --+ (.5,0) circle [radius=.05cm];
	
	\draw[thick, draw=teal, fill=teal] ($(huvudmodul.south east) + (0,.25)$) --+ (.25,0) circle [radius=.05cm];
	\draw[thick, draw=teal, fill=teal] ($(huvudmodul.south east) + (0,.5)$) --+ (.5,0) circle [radius=.05cm];
	
	\draw[thick, draw=teal, fill=teal] ($(styrmodul.south east) + (0,.25)$) --+ (.25,0) circle [radius=.05cm];
	\draw[thick, draw=teal, fill=teal] ($(styrmodul.south east) + (0,.5)$) --+ (.5,0) circle [radius=.05cm];
	
	%Koppling till sensormodul
	\draw[thick, draw=violet] (.5,9.5) --++ (.25,0) --++ (0,-2.5) --++ (0.75,0) ;
	\draw[thick, draw=violet] (.5,.5) --++ (.25,0) --++ (0,6.25) --++ (0.75,0);
	
	\draw[thick, draw=violet] (5.5,9.5) --++ (-.25,0) --++ (0,-2.5) --++ (-0.75,0);
	\draw[thick, draw=violet] (5.5,.5) --++ (-.25,0) --++ (0,6.25) --++ (-0.75,0);
	
	\draw[thick, draw=violet] (2,9.5) --++ (0,-2);
	\draw[thick, draw=violet] (4,9.5) --++ (0,-2);
	
	\draw[thick, draw=violet] (sensormodul.south) --++ (0,-.5);
	
	%Kopplingar till styrmodul
	\draw[thick, draw=olive] (styrmodul.east) --++ (1.5,0);
	\draw[thick, draw=olive] (styrmodul.west) --++ (-1.5,0);
	\draw[thick, draw=olive] ($(styrmodul.east) + (0,.25)$) --++ (1,0) --++ (0,6.25) --++ (.5,0);
	\draw[thick, draw=olive] ($(styrmodul.west) + (0,.25)$) --++ (-1,0) --++ (0,6.25) --++ (-.5,0);
	\draw[thick, draw=olive] (3,9.5) --++ (0,-1.75) --++ (-1.75,0) --++ (0,-6) --++ (.25,0);
	
	%Bluetooth
	\draw[thick, draw=black] (huvudmodul.east) --++ (2,0) node[right,minimum width=.75cm, minimum height=1cm, draw=black] (bt) {\includegraphics{bluetooth}};
	
	\node (bt2) at ($(bt) + (2,0)$) [thick, draw=black, minimum width=.75cm, minimum height=1cm] {\includegraphics{bluetooth}};
	
	\draw[thick, ->,line join=round,decorate, decoration={
    												snake,
    												segment length=5,
    												amplitude=1,
    												post=lineto,
    												post length=1pt}] 
    		($(bt.east) + (5pt,5pt)$) -- ($(bt2.west) + (-5pt,5pt)$);
    		
    \draw[thick, ->,line join=round,decorate, decoration={
    												snake,
    												segment length=5,
    												amplitude=1,
    												post=lineto,
    												post length=1pt}] 
    		 ($(bt2.west) + (-5pt,-5pt)$) -- ($(bt.east) + (5pt,-5pt)$);
    		 
    \draw[thick, draw=black] (bt2.east) --++ (.5,0) node[right, minimum width=3cm, minimum height=1.5cm, draw=black, fill=white] {Datormodul};
    
    \draw[thick, draw=black] ($(sensormodul.south west) + (.25,0)$) --++ (0,-2) node[midway, above, sloped] {\small avbrott};
    
    \draw[thick, draw=black] ($(styrmodul.north west) + (.25,0)$) -- ($(huvudmodul.south west) + (.25,0)$) node[right, midway] {\small avbrott};
	
	\end{tikzpicture}
	
\end{document}  
}
	\caption{Det totala systemet \label{overview}}	
\end{figure}

Med hjälp av en avsökningsalgoritm kan roboten navigera i ett labyrintsystem. Denna är beroede av sensorer för att få en uppfattning om sin omgivning samt en styrenhet för att kunna röra sig i förhållande till sin omgivning. Navigationen pågår tills den detekterar en nödställd. Då ser den till att finna den kortaste vägen ut för att där hämta en förnödenhet med gripklon som den sedan åker och lämnar till den nödställde. 

\section{Teori}
\textit{Beskrivning av regleralgoritmer mm.}

\section{Systemet}
\textit{Ett översiktligt blockschema och en beskrivning av hela systemet.}

\subsection{Protokoll}

\section{Modulerna}
\textit{Innehåller mera detaljerade blockschemor och beskrivningar av varje modul. Tänk på läsbarheten och växla mellan figurer och text.}
I detta avsnitt följer en detaljerad beskrivning av systemets ingående moduler. 

\subsection{Huvudmodulen}
Huvudmodulen har som uppgift att ta emot och förmedla information mellan de andra modulera samt tar alla övergripande beslut. Bluetooth\textsuperscript{\circledR} används för att kommunicera med datormodulen och en I\textsuperscript{2}C-buss används för att kommunicera med sensor- och styrmodulen, vilket visas i figur \ref{communication}.

\begin{figure}[htbp]
\noindent\resizebox{.97\textwidth}{!}{
	\documentclass[border=20pt]{standalone}
\usepackage{tikz}
\usetikzlibrary{positioning}
\usetikzlibrary{calc}
\usetikzlibrary{decorations.pathmorphing}
\usepackage{amssymb}
\usetikzlibrary{shapes,arrows}

\begin{document}
	\begin{tikzpicture}[scale=1]
		
		\tikzset{every node/.style={thick, draw=black, align=center, minimum height=40pt, text width=100pt, minimum width=100pt}}
		\node(datormodul) {Datormodul};
		\node[right=10pt of datormodul,minimum height=20pt, minimum width=10pt,text width=10pt] (bt1) {\includegraphics{bluetooth}};
		
		\node[right=40pt of bt1,minimum height=20pt, minimum width=10pt,text width=10pt] (bt2) {\includegraphics{bluetooth}};
		
		\node[right=10pt of bt2] 			(huvudmodul) 	{Huvudmodul};
		\node[below=-10pt of huvudmodul,draw=none] (master) {\textit{master}};
		\node[right=10pt of huvudmodul] 		(sensormodul) 	{Sensormodul};
		\node[below=-10pt of sensormodul,draw=none] (slave1) {\textit{slav}};
		\node[right=10pt of sensormodul] 	(styrmodul) 		{Styrmodul};
		\node[below=-10pt of styrmodul,draw=none] (slave2) {\textit{slav}};
		
		\coordinate (sclStart) 	at ($(huvudmodul.north west) + (0,20pt)$);
		\coordinate (sclEnd)		at ($(styrmodul.north east)  + (0,20pt)$);
		
		\coordinate (sdaStart)  at ($(sclStart) + (0,20pt)$);
		\coordinate (sdaEnd)		at ($(sclEnd)	+ (0,20pt)$);
		
		\coordinate (vddStart)  at ($(sdaStart) + (0,40pt)$);
		\coordinate (vddEnd)		at ($(sdaEnd)	+ (0,40pt)$);
		
		\draw[thick] (sclStart) -- (sclEnd) node [right,draw=none,text width=0,minimum width=0] {SCL};
		\draw[thick] (sdaStart) -- (sdaEnd) node [right,draw=none,text width=0,minimum width=0] {SDA};
		\draw[thick] (vddStart) -- (vddEnd) node [right,draw=none,text width=0,minimum width=0] {$V_{dd}$};
		
		\draw[thick] (datormodul.east) -- (bt1.west);
		
		\draw[thick, ->,line join=round,decorate, decoration={
    												snake,
    												segment length=5,
    												amplitude=1,
    												post=lineto,
    												post length=1pt}] 
    		($(bt1.east) + (5pt,5pt)$) -- ($(bt2.west) + (-5pt,5pt)$);
    		
    	\draw[thick, ->,line join=round,decorate, decoration={
    												snake,
    												segment length=5,
    												amplitude=1,
    												post=lineto,
    												post length=1pt}] 
    		 ($(bt2.west) + (-5pt,-5pt)$) -- ($(bt1.east) + (5pt,-5pt)$);
    		 
    	\draw[thick] (bt2.east) -- (huvudmodul.west);
    	
    	\draw[thick,fill=black] ($(huvudmodul.north) + (-10pt,0)$) -- ($(huvudmodul.north) + (-10pt,20pt)$) circle [radius=2pt];
    	\draw[thick,fill=black] ($(huvudmodul.north) + (+10pt,0)$) -- ($(huvudmodul.north) + (+10pt,40pt)$) circle [radius=2pt];
    	
    	\draw[thick,fill=black] ($(sensormodul.north) + (-10pt,0)$) -- ($(sensormodul.north) + (-10pt,20pt)$) circle [radius=2pt];
    	\draw[thick,fill=black] ($(sensormodul.north) + (+10pt,0)$) -- ($(sensormodul.north) + (+10pt,40pt)$) circle [radius=2pt];
    	
    	\draw[thick,fill=black] ($(styrmodul.north) + (-10pt,0)$) -- ($(styrmodul.north) + (-10pt,20pt)$) circle [radius=2pt];
    	\draw[thick,fill=black] ($(styrmodul.north) + (+10pt,0)$) -- ($(styrmodul.north) + (+10pt,40pt)$) circle [radius=2pt];
    	
    	\draw[thick,fill=black] ($(sensormodul.north east) + (-5pt,20pt)$) circle [radius=2pt] -- ($(sensormodul.north east) + (-5pt,80pt)$) circle [radius=2pt];
    	\draw[thick,fill=black] ($(styrmodul.north west) + (5pt,40pt)$) circle [radius=2pt] -- ($(styrmodul.north west) + (5pt,80pt)$) circle [radius=2pt];
    	\draw[thick,draw=black,fill=white] ($(styrmodul.north west) + (2pt,50pt)$) rectangle ($(styrmodul.north west) + (8pt,70pt)$);
		\draw[thick,draw=black,fill=white] ($(sensormodul.north east) + (-8pt,50pt)$) rectangle ($(sensormodul.north east) + (-2pt,70pt)$) node [below left=-10pt and 20pt,draw=none,minimum width=0, text width = 0pt] {$R_p$};
		
		\draw[thick, ->] ($(sensormodul.south) + (+40pt,0)$) -- ($(sensormodul.south) + (40pt,-30pt)$) -- ($(huvudmodul.south) + (40pt,-30pt)$) -- ($(huvudmodul.south) + (40pt,0)$);
		\draw[thick, ->] ($(styrmodul.south) + (+40pt,0)$) -- ($(styrmodul.south) + (40pt,-40pt)$) -- node [midway,below,minimum height=20pt, draw=none] {Avbrott}($(huvudmodul.south) + (30pt,-40pt)$) -- ($(huvudmodul.south) + (30pt,0)$);
	\end{tikzpicture}
\end{document}}
	\caption{Intermodulär kommunikation \label{communication}}
\end{figure}

\pagebreak

\subsection{Sensormodulen}
Robotens sensormodul har som uppgift att läsa av sensorerna och kommunicera datan till huvudmodulen. De sensorer som används är följande, NAMN PÅ SENSORER
\begin{description}
	\item[Avstånd] \hfill \\
	4 x IR-sensor \verb+GP2D120+ (4 cm till 30 cm) \\
	1 x Laser-sensor \verb+LIDAR-Lite v2+ (0 till 40 m) \\
	\item[Vinkelhastighet] \hfill \\
	1 x Gyro/accelerometer \verb+MPU-6500+ 
	\item[Identifierare av nödställd] \hfill \\
	1 x IR-detektor \verb+IRM-8601-S+
\end{description}
och är placerade enligt figur \ref{sensors}. 

\begin{figure}[htbp]
\centering
\noindent\resizebox{.8\textwidth}{!}{
	\documentclass[border=10px]{standalone}
\usepackage{tikz}
\usetikzlibrary{patterns}
\usetikzlibrary{shapes.arrows}
\usepackage{amssymb}
\usetikzlibrary{calc}
\usepackage{verbatim}
\begin{document}
	
\begin{tikzpicture}[scale=1,rotate=90]
		
	%Base
	\draw[thick, draw=black, fill=gray!10] (0,0) rectangle (6,10);

	%Wheels
	\draw[thick, pattern=north west lines, pattern color=black] (-.5,1) 		rectangle (0,2.5);
	\draw[thick, pattern=north west lines, pattern color=black] (-.5,7.5) 	rectangle (0,9);
	\draw[thick, pattern=north west lines, pattern color=black] (6,1) 		rectangle (6.5,2.5);
	\draw[thick, pattern=north west lines, pattern color=black] (6,7.5) 		rectangle (6.5,9);
	
	%Sensors
	\draw[thick, draw=black, fill=white] (-.25,.25) 		rectangle (.5,.75);
	\draw[thick, draw=black, fill=white] (-.25,9.25) 	rectangle (.5,9.75);
	\draw[thick, draw=black, fill=white] (5.5,.25) 		rectangle (6.25,.75);
	\draw[thick, draw=black, fill=white] (5.5,9.25) 		rectangle (6.25,9.75);
	\draw[thick, draw=black, fill=white] (2,10.25) 		rectangle (4,9.5);
	\draw[thick, draw=black, fill=white] (2.5,4) 		rectangle (3.5,6);
	
	%Arrows and text
	\draw[thick, ->]  (3,11) node[left, align=center] {\verb+LIDAR-Lite v2+ \\ + detektor av nödställd} -- (3,10.25);
	\draw[thick, <->] (0.5,0.5)  --  (5.5,0.5) node[left=-14pt,midway, fill=gray!10] {\verb+GP2D120+};
	\draw[thick, <->] (0.5,9.5) -- (3,9) node[right=-14pt,fill=gray!10] {\verb+GP2D120+} -- (5.5,9.5);
	\draw[thick, ->] (4.5,5) node[above] {\verb+MLX90609+} -- (3.5,5);
	\end{tikzpicture}
	
\end{document}	}
	\caption{Sensorplacering \label{sensors}}
\end{figure}

\subsection{Styrmodulen}
Styrmodulen är den del av roboten som tar emot och verkställer styrkommandon och ser till att roboten tar sig fram på önskat sätt. Det som ska kunna hanteras är styrning av chassits motorer för att få roboten att röra på sig, rotera det eventuella lasersensortornet och att styra gripklon.

\subsection{Datormodulen}

\section{Slutsatser}
\textit{Vilka förbättringar skulle kunna göras?}

\pagebreak
\addcontentsline{toc}{section}{Referenser}
\bibliographystyle{ieeetr}
\bibliography{references}

\pagebreak

\appendix

\end{flushleft}

\end{document}